% Generated deterministically from the audited Markdown source.
% Responsible author: Carptopus.
\documentclass[11pt]{article}
\usepackage[a4paper,margin=27mm]{geometry}
\usepackage[T1]{fontenc}
\usepackage[utf8]{inputenc}
\usepackage{lmodern}
\usepackage{microtype}
\usepackage{amsmath,amssymb,amsthm,mathtools}
\usepackage{booktabs,array}
\usepackage{xcolor}
\usepackage[colorlinks=true,linkcolor=blue!55!black,citecolor=blue!55!black,urlcolor=blue!65!black]{hyperref}
\usepackage{xurl}
\usepackage{fancyhdr}

\newtheorem{theorem}{Theorem}
\newtheorem{lemma}[theorem]{Lemma}
\newtheorem{proposition}[theorem]{Proposition}
\pagestyle{fancy}
\fancyhf{}
\fancyhead[L]{Carptopus}
\fancyhead[R]{Odd-dimensional RM2 descent obstructions}
\fancyfoot[C]{\thepage}
\setlength{\headheight}{14pt}
\setlength{\parindent}{0pt}
\setlength{\parskip}{0.42em}
\setlength{\emergencystretch}{4em}
\urlstyle{same}

\title{Odd-dimensional descent obstructions\\
in the third support spectrum of binary second-order Reed--Muller codes}
\author{Carptopus\\\href{mailto:carptopus@163.com}{carptopus@163.com}}
\date{Manuscript, 2 September 2026}

\hypersetup{
  pdftitle={Odd-dimensional descent obstructions in the third support spectrum of binary second-order Reed-Muller codes},
  pdfauthor={Carptopus},
  pdfsubject={A local parity band and infinite descent obstructions in the third RM2 support spectrum},
  pdfkeywords={Reed-Muller codes, higher support spectrum, quadratic Boolean functions, Walsh transform, Fano plane, Pfaffian}
}

\begin{document}
\maketitle

\begin{center}
\fcolorbox{black!35}{black!4}{\parbox{0.88\textwidth}{\centering\small
Internally verified candidate manuscript, version 0.1-beta. The mathematical claims and the
complete manuscript have passed project-internal zero-trust audits; external review is pending.}}
\end{center}

\begin{abstract}


Let $\mathcal S_n^{(3)}$ be the support spectrum of three-dimensional subcodes of the binary
second-order Reed--Muller code $RM_2(2,n)$. We prove that the tempting descent relation
$\mathcal S_{2m+1}^{(3)}=2\mathcal S_{2m}^{(3)}$ fails for every $m\geq5$, by an explicit
infinite family. More generally, after normalizing the seven zero-frequency Walsh coefficients of
a quadratic net, we classify an exponentially widening band around the aggregate value
$3\cdot2^{m-1}$. For every $m\geq6$ and every odd $c$ with $|c|<2^{m-4}$, the value
$3\cdot2^{m-1}+c$ is attainable in odd dimension if and only if
$c\in\{\pm1,\pm3,\pm5\}$, and no value in the same band is attainable in even dimension.
The proof combines explicit quadratic nets with Fano-plane half-rank propagation and scalar and
vector maximal-Pfaffian identities. This is a local parity classification, not a determination of
the complete third support spectrum.

\end{abstract}

\section{Introduction}


For a binary linear code $C\subseteq\mathbb F_2^N$ and a subcode $D\leq C$, the support of $D$
is the union of the supports of its codewords. Put

\[
\mathcal S_n^{(3)}=
\{|\operatorname{Supp}(D)|:D\leq RM_2(2,n),\ \dim D=3\}.
\]

Earlier work determined several exact regions of this spectrum, including the values for
$2\leq n\leq9$, two sharp all-dimensional gaps, a near-full interval, and a low-polar-rank
recursion \cite{rm2regions}. A second companion paper classifies the labelled Fano half-rank profiles of
three-parameter alternating nets and their minimum dimensions \cite{fanoprofiles}; for the wider higher-weight
spectrum context, see \cite{higherweights}. The systems isolated below all contain nonzero-Walsh polar-rank-two
directions and therefore lie inside the complete finite recursion of \cite[Theorem 6.1]{rm2regions}. The new
increment is not a previously uncovered stratum: it is a closed-form all-parameter extraction from
that recursion, together with a direct proof, an exact local parity band and an infinite descent
obstruction. Neither earlier paper states these conclusions.

The first new phenomenon is an infinite family of odd-dimensional supports whose halves are
forbidden one dimension lower. The stronger result shows that this is not an isolated ray: in a
band whose width grows exponentially with dimension, exactly six odd-dimensional offsets survive
once $m\geq7$, whereas the entire even-dimensional band is empty. At $m=6$, the band itself
contains four of them and a separate endpoint argument supplies the remaining two.

\section{Walsh aggregates and support sizes}


Let $D=\langle f_1,f_2,f_3\rangle$ be a three-dimensional quadratic subcode and write

\[
f_a=a_1f_1+a_2f_2+a_3f_3,
\qquad a\in\mathbb F_2^3\setminus\{0\}.
\]

For a Boolean function $f$ on $\mathbb F_2^n$, let

\[
W_f(0)=\sum_{x\in\mathbb F_2^n}(-1)^{f(x)}.
\]

Character orthogonality gives

\[
|\operatorname{Supp}(D)|
=7\cdot2^{n-3}-\frac18\sum_{a\ne0}W_{f_a}(0).
\tag{1}
\]

For even and odd dimensions define

\[
\begin{aligned}
\mathcal K_m^{\rm ev}
&=\left\{\sum_{a\ne0}2^{-m}W_{f_a}(0):
 D\leq RM_2(2,2m),\ \dim D=3\right\},\\
\mathcal K_m^{\rm odd}
&=\left\{\sum_{a\ne0}2^{-(m+1)}W_{f_a}(0):
 D\leq RM_2(2,2m+1),\ \dim D=3\right\}.
\end{aligned}
\tag{2}
\]

Thus the same normalized aggregate $K$ gives the odd support

\[
7\cdot2^{2m-2}-2^{m-2}K
\tag{3}
\]

and, if attainable, the even support

\[
7\cdot2^{2m-3}-2^{m-3}K,
\tag{4}
\]

which is exactly half of (3).

Let $B_a$ be the polar form of $f_a$ and put
$h_a=\operatorname{rank}(B_a)/2$. On every Fano line $\{a,b,a+b\}$,

\[
|h_a-h_b|\leq h_{a+b}\leq h_a+h_b.
\tag{5}
\]

When equality holds on the right and all three Walsh sums are nonzero, quadratic Gauss-sum
multiplication fixes their relative signs. In the normalization (2), every coordinate lies in

\[
E_m=\{0,-2^m\}\cup\{\pm2^j:0\leq j<m\}.
\tag{6}
\]

The value $+2^m$ would come from the zero Boolean function. It cannot occur at a nonzero parameter
$a$ because $f_1,f_2,f_3$ are linearly independent; the constant-one function accounts for the
possible value $-2^m$.

These elementary restrictions are the input to the classification below.

\section{The local parity band}


Put $A=2^{m-1}$.

\begin{theorem}


For every $m\geq6$ and every odd integer $c$ satisfying $|c|<A/8$,

\[
3A+c\in\mathcal K_m^{\rm odd}
\quad\Longleftrightarrow\quad
c\in\{-5,-3,-1,1,3,5\},
\tag{7}
\]

whereas

\[
3A+c\notin\mathcal K_m^{\rm ev}
\tag{8}
\]

for every such $c$.

We first reduce all possible seven-coordinate patterns.

\end{theorem}

\begin{lemma}


Suppose seven values in $E_m$ have sum $3A+c$, satisfy (5) and the equality-sign rule, and
$m\geq6$, $c$ is odd, and $|c|<A/8$. Then exactly three coordinates equal $+A$; they form a
Fano line. The other four coordinates lie in $\{0,\pm1,\pm2\}$, and consequently
$c\in\{\pm1,\pm3,\pm5,\pm7\}$.

\end{lemma}

\begin{proof}


Let $q$ be the number of coordinates equal to $+A$.

If $q=0$, let $r$ count the $+A/2$ coordinates. For $r\leq4$ the total is at most
$11A/4<3A+c$. For $r=5$, the remaining two coordinates can reach the required total only at the
excluded central value $c=0$; dropping either coordinate to the next binary layer loses at least
$A/8$. If $r=6$, the last coordinate would equal the odd number $c$, so only $c=\pm1$ is
compatible with (6). Its half-rank is $m$, while the other six have half-rank two; expressing the
exceptional Fano point as a sum of two of the other points contradicts (5). The case $r=7$ is
immediate.

If $q=1$, call the unique point $p$. Translation by $p$ partitions the remaining six points into
three pairs. Their total is the odd number $2A+c$, so one pair contains a coordinate of absolute
value one. That pair contributes at most three, while each other pair contributes at most $A$.
If $c>3$, this upper bound is already smaller than $2A+c$. If $c\leq3$, the hypothesis
$|c|<A/8$ forces the four points in the two high pairs to equal $+A/2$;
otherwise one binary-layer drop is already too large. The two cross-sums of these points form the
third pair and have half-rank at most four, contradicting the full-rank coordinate in that pair.

If $q=2$, the sum of the two $+A$ points has half-rank at most two and contributes at most $A/2$.
The complementary affine plane contains a full-rank point, and repeated use of (5) bounds the sum
of the absolute values on that plane by nine. But $A/2+9<A+c$ for $A\geq32$, a contradiction.

If $q\geq4$, four rank-two points and their pairwise sums cover the Fano plane. Hence every
half-rank is at most two and the total is divisible by $A/2$, contrary to the odd offset $c$.

It remains that $q=3$. Any additional coordinate of half-rank at most one would again propagate
to four low-rank points and force the preceding divisibility contradiction. If the three $+A$
points are not collinear, all four remaining values are multiples of $A/4$, again impossible for
odd $c$. They therefore form a line. Any two points in the complementary affine plane sum to a
point of this line, so their half-ranks differ by at most one. Since their sum $c$ is odd, one is
full rank; all four half-ranks are consequently $m-1$ or $m$, and their values belong to
$\{0,\pm1,\pm2\}$. Their odd sum has absolute value at most eight, giving the asserted eight
possibilities.
\end{proof}


Two Pfaffian identities eliminate the two extreme offsets and separate the parities. In even
dimension, if $R,S,R+S$ all have rank at most two, the maximal Pfaffian restricted to the affine
plane $C+\langle R,S\rangle$ is affine linear. It is therefore nonzero at zero, two, or four
points. In odd dimension $2m+1$, define the maximal-Pfaffian vector

\[
P(C)=\bigl(\operatorname{Pf}(C_{\widehat i})\bigr)_{i=1}^{2m+1}.
\]

The same rank-two update argument, applied coordinatewise, gives

\[
P(C+R)+P(C+S)+P(C+R+S)=P(C).
\tag{9}
\]

Moreover $P(C)\ne0$ exactly when $C$ has rank $2m$. Hence an affine plane cannot contain exactly
one maximal-rank form.

In even dimension a full-rank quadratic function has nonzero Walsh sum. The number of full-rank
coordinates among the complementary four points is therefore the number of entries of absolute
value one, which has the same parity as odd $c$. The scalar Pfaffian restriction says this number
must be zero, two, or four, a contradiction. This proves (8).

In odd dimension, the offsets $c=7$ and $c=-7$ force respectively the multisets
$\{1,2,2,2\}$ and $\{-2,-2,-2,-1\}$. Each has exactly one maximal-rank polar form, contradicting
(9). Lemma 2 has already excluded every other odd offset outside
$\{\pm1,\pm3,\pm5\}$. It remains only to construct these six values.

\section{Six explicit odd-dimensional families}


On five variables define

\[
\begin{aligned}
q_1={}&x_0x_3+x_1x_2+x_1x_3+x_2x_3+x_2x_4+x_3x_4,\\
q_2={}&x_0x_1+x_0x_2+x_0x_3+x_0x_4+x_1x_4+x_2x_4+x_3x_4,\\
q_3={}&x_0x_1+x_0x_2+x_0x_3+x_0x_4+x_1x_2+x_1x_3+x_1x_4.
\end{aligned}
\tag{10}
\]

Encode an affine function on these variables by a six-bit mask, with the constant term first.
The following masks give the listed normalized seven-coordinate Walsh signatures. Coordinates are
ordered by the binary parameters $a=1,2,\ldots,7$, with the least significant bit multiplying the
first generator.

\begin{table}[ht]
\centering
\small
\begin{tabular}{r l l}
\toprule
$c$ & $(\ell_1,\ell_2,\ell_3)$ & base signature \\
\midrule
$-5$ & $(1,2,2)$ & $(-1,2,-1,2,-1,2,-2)$ \\
$-3$ & $(1,2,63)$ & $(-1,2,-1,2,-1,2,0)$ \\
$-1$ & $(1,28,4)$ & $(-1,2,0,2,0,2,0)$ \\
$1$ & $(4,2,63)$ & $(0,2,0,2,1,2,0)$ \\
$3$ & $(0,2,63)$ & $(1,2,1,2,1,2,0)$ \\
$5$ & $(0,2,2)$ & $(1,2,1,2,1,2,2)$ \\
\bottomrule
\end{tabular}
\caption{Affine masks and normalized base signatures for the six explicit families.}
\end{table}


For $m\geq3$, set

\[
\begin{aligned}
F_{m,c,1}&=q_1+\ell_1+\sum_{j=0}^{m-3}x_{5+2j}x_{6+2j},\\
F_{m,c,2}&=q_2+\ell_2,\\
F_{m,c,3}&=q_3+\ell_3.
\end{aligned}
\tag{11}
\]

Each added hyperbolic pair multiplies the seven normalized coordinates by
$(1,2,1,2,1,2,1)$. Hence the three even-position coordinates grow from two to $A$, while the
other four retain sum $c$. Thus

\[
\sum_{a\ne0}2^{-(m+1)}W_{F_{m,c,a}}(0)=3A+c.
\tag{12}
\]

Here $F_{m,c,a}=a_1F_{m,c,1}+a_2F_{m,c,2}+a_3F_{m,c,3}$.

The three polar components are linearly independent, so these are genuine three-dimensional
subcodes. Equations (7) and (8) now follow from (12) and the preceding exclusions, completing the
proof of Theorem 1.

\section{Endpoint completion and descent failure}


For $c\in\{\pm1,\pm3,\pm5\}$ define

\[
s_{m,c}=2^{m-2}(11A-c).
\tag{13}
\]

For $m\geq7$, all six offsets satisfy $|c|<A/8$, so Theorem 1 immediately gives the following
six-family statement. The excluded endpoints at $m=6$ require a short boundary check.

\begin{proposition}


For every $m\geq6$ and $c\in\{\pm1,\pm3,\pm5\}$,

\[
s_{m,c}\in\mathcal S_{2m+1}^{(3)},
\qquad
s_{m,c}/2\notin\mathcal S_{2m}^{(3)}.
\tag{14}
\]

\end{proposition}

\begin{proof}


Theorem 1 and the constructions in Section 4 prove the claim for $m\geq7$ and for
$m=6$, $c=\pm1,\pm3$. It remains to exclude $K=3A\pm5$ in even dimension when
$m=6$, so $A=32$; odd attainability is already given by (11).

Repeat the $q$-partition from Lemma 2. For $q=0$, if at most four coordinates equal $+16$, the
total is below $91$. With five such coordinates, the last two would have to sum to $21$ or $11$;
the first exceeds their maximum $16$, while the second is not a sum of two signed powers of two
from $\{0,\pm1,\pm2,\pm4,\pm8\}$. With six such coordinates the last value would be $\pm5$,
which is not in $E_6$, and seven give the wrong total.

For $q=1$, the full-rank pair contributes at most three. The offset $+5$ exceeds the resulting
upper bound. For the offset $-5$, equality in the only remaining layer budget would force three
of the four high coordinates to be $+16$, the fourth to be $+8$, and the full-rank pair to be
$\{1,2\}$. The four high half-ranks are $2,2,2,3$; each point of the remaining pair is a sum of
two of them and hence has half-rank at most five, contradicting the required half-rank six.
The $q=2$, $q\geq4$, and $q=3$ arguments of Lemma 2 use only the inequalities
$A/2+9<A-5$, divisibility, and Pfaffian parity, respectively, and apply unchanged. Thus both
endpoints are even-dimensionally impossible.
\end{proof}


The offset $c=5$ begins one dimension pair earlier. We record the missing endpoint explicitly.

\begin{proposition}


For every $m\geq5$, with

\[
s_m=2^{m-2}(11\cdot2^{m-1}-5),
\]

we obtain

\[
s_m\in\mathcal S_{2m+1}^{(3)},
\qquad
s_m/2\notin\mathcal S_{2m}^{(3)}
\quad(m\geq5).
\tag{15}
\]

\end{proposition}

\begin{proof}


For $m\geq6$ this is Proposition 3. The construction (11) for $c=5$ is valid already at
$m=5$, so only even-dimensional exclusion remains. Here $A=16$ and the target aggregate is 53.

Let $q$ again count the $+A$ coordinates. If $q=0$, at most five $+8$ coordinates give total at
most $5\cdot8+2\cdot4=48$; six would require the last coordinate to equal five, and seven sum to
56. If $q=1$, the three translated pairs contribute at most $3+16+16=35$ beyond the unique
$16$, below 53. If $q=2$, the sum point contributes at most eight and the complementary affine
plane at most nine in absolute value, below the required residual 21. If $q\geq4$, the total is a
multiple of eight. Hence $q=3$.

The three $+16$ points must form a Fano line. The equality-sign rule and the half-rank triangle
inequalities leave only the complementary multisets $\{0,1,2,2\}$ and
$\{1,1,1,2\}$. They contain respectively one and three full-rank forms. The scalar maximal
Pfaffian on that affine plane is affine linear, so the number of full-rank points must be zero, two,
or four. Both patterns are impossible. This proves the $m=5$ exclusion.
\end{proof}


In particular $\mathcal S_{2m+1}^{(3)}\ne2\mathcal S_{2m}^{(3)}$ for every $m\geq5$.

\section{Scope and verification}


Theorem 1 is exact only inside the band $|c|<2^{m-4}$ around $3\cdot2^{m-1}$. It does not
classify the full difference $\mathcal K_m^{\rm odd}\setminus\mathcal K_m^{\rm ev}$, determine
the complete support spectrum, or give a finite recursion for all quadratic nets. The bandwidth is
a clean sufficient range for the layer-gap argument and is not claimed optimal.

The source Loop directories contain exact regression checks for the explicit signatures, support
formula and finite parameter compression. In particular,
{\small\path{verify_six_endpoint_families.py}} reuses the Walsh-transform helpers developed during the earlier
RM2 investigation, while {\small\path{verify_local_slope_three_band.py}} checks the finite compressed cases used
to discover and calibrate the local band. These are not independent implementations of the whole
argument, and the Pfaffian exclusions are proved in the manuscript rather than certified by those
scripts. The all-$m$ conclusions therefore depend on the proofs above, not on bounded enumeration.

\begin{thebibliography}{9}
\footnotesize
\raggedright

\bibitem{rm2regions}
Carptopus,
\emph{Exact regions and low-polar-rank recursion in the third support spectrum of binary
second-order Reed--Muller codes}, Public Beta v0.1 (2026).
\href{https://doi.org/10.5281/zenodo.22165319}{doi:10.5281/zenodo.22165319}.

\bibitem{fanoprofiles}
Carptopus,
\emph{Realizability and minimum dimension of Fano half-rank profiles over $\mathbb F_2$},
Public Beta v0.1 (2026).
\href{https://doi.org/10.5281/zenodo.22164877}{doi:10.5281/zenodo.22164877}.

\bibitem{higherweights}
S.~R. Ghorpade, T. Johnsen, A.~R. Ludhani and P. Pratihar,
\emph{Higher weight spectra of codes},
\href{https://arxiv.org/abs/2408.02548}{arXiv:2408.02548}.

\end{thebibliography}

\paragraph{AI-assisted research and writing disclosure.}
Carptopus is the author responsible for the mathematical claims and the public version of this
manuscript. OpenAI Codex was used as an AI-assisted tool for literature organization, proof
development and stress testing, exact verification, adversarial auditing, and manuscript
preparation.

\end{document}
